% !TeX root = ../main.tex
\chapter{Theoretical and Technological Foundations}
\label{chap:co-so-ly-thuyet-cong-nghe}

This chapter presents the core theoretical and technical foundations used
in the LMS system that automatically generates Micro-Content and Quizzes
with Generative AI. The chapter focuses on four main groups of content:
Micro-learning and Bloom Taxonomy, Large Language Models,
Retrieval-Augmented Generation, and chunking/document-processing
techniques. These foundations support the design of the document
processing, content generation, quiz generation, and learning chatbot
pipelines in later chapters.

\section{Micro-learning and Bloom Taxonomy}
\label{sec:microlearning-bloom}

\subsection{Micro-learning}
\label{subsec:microlearning}

Micro-learning organizes learning content into small units, each focused
on a specific concept, skill, or learning objective. Instead of requiring
learners to absorb a long chapter or a lecture video lasting tens of
minutes, micro-learning divides content into short units such as lesson
summaries, key concept cards, flashcards, self-check questions, or short
explanations with examples.

In the LMS context, micro-learning has three important meanings. First,
it helps learners approach knowledge in short learning rhythms, suitable
for mobile learning habits and short intermittent study sessions. Second,
it allows the system to track progress at a finer granularity: which card
the learner has read, which concept they answered incorrectly, or which
Learning Outcome should be reviewed. Third, it provides a clear output
format for AI generation: from one document chunk, the system can
generate a structured set of cards and quizzes.

For this thesis, micro-learning is not only a presentation interface but
also a direct influence on the document-processing pipeline. Input
documents must be extracted, chunked, attached with source metadata, and
then transformed into short learning units that can be reviewed and
published. Each micro-content unit must maintain its relationship with
the original document, its position in the chapter, the related Learning
Outcome, and the instructor review state.

\begin{table}[H]
    \centering
    \caption{Characteristics of micro-learning in the proposed system}
    \label{tab:microlearning-characteristics}
    \begin{tabular}{p{4cm}p{9.5cm}}
        \toprule
        \textbf{Characteristic} & \textbf{Meaning in the system} \\
        \midrule
        Concise & Each content unit focuses on one main idea, reducing cognitive load for learners. \\
        Source-aware & Cards/quizzes must link to a document, page, heading, or video timestamp for verification. \\
        Quickly assessable & Each content cluster can attach flashcards or quizzes to reinforce memory. \\
        Pedagogical metadata & Content is tagged with Learning Outcome, Bloom level, difficulty, and review state. \\
        Personalization-friendly & Completion data for each card/quiz can be used to recommend review content. \\
        \bottomrule
    \end{tabular}
\end{table}

\subsection{Bloom Taxonomy}
\label{subsec:bloom-taxonomy}

Bloom Taxonomy is a framework for classifying learning objectives by
cognitive level. In the commonly used revised classification, cognitive
levels include Remember, Understand, Apply, Analyze, Evaluate, and
Create. These levels represent increasing cognitive complexity, from
remembering basic information to evaluating and creating new products.

In quiz generation, Bloom Taxonomy is used as a pedagogical metadata
layer to guide question design. If the system only generates questions
from surface content, quizzes tend to remain at the simple memorization
level. With Bloom level metadata, instructors can ask the system to
create questions for specific objectives, for example testing the ability
to ``explain'', ``apply'', or ``analyze'' instead of merely ``list''.

\begin{longtable}{|p{2.6cm}|p{4.2cm}|p{7cm}|}
\caption{Bloom levels and their meaning in quiz generation}
\label{tab:bloom-levels}\\
\hline
\textbf{Bloom level} & \textbf{Common verbs} & \textbf{Meaning in question generation} \\
\hline
\endfirsthead
\hline
\textbf{Bloom level} & \textbf{Common verbs} & \textbf{Meaning in question generation} \\
\hline
\endhead
Remember & list, define, identify, name & Tests the ability to remember terms, concepts, formulas, or facts. \\
\hline
Understand & explain, describe, distinguish, summarize & Tests the ability to interpret meaning and recognize basic relations among concepts. \\
\hline
Apply & apply, calculate, use, illustrate & Tests the ability to use knowledge in a concrete situation or example. \\
\hline
Analyze & analyze, compare, infer, classify & Tests the ability to separate a problem into components and recognize causes, consequences, or structure. \\
\hline
Evaluate & evaluate, select, critique, justify & Tests the ability to make judgments based on criteria. This level is often more suitable for essays than simple MCQs. \\
\hline
Create & design, build, propose, create & Tests the ability to produce a new product or solution. This level usually requires projects or open-ended questions. \\
\hline
\end{longtable}

Within this thesis, Remember, Understand, Apply, and Analyze are the most
suitable levels for automatic multiple-choice quizzes. Evaluate and
Create are still stored in the metadata model for future expansion, but
they are not the focus of the MVP because they usually require essay
grading, rubrics, or project-based assignments.

\subsection{Bloom-aware Quiz Generation}
\label{subsec:bloom-aware-quiz-generation}

Bloom-aware Quiz Generation is the process of generating questions with
explicit awareness of cognitive level. Instead of simply asking the LLM
to ``create five questions from the passage'', the system provides
constraints about Learning Outcome, Bloom level, difficulty, question
type, and output structure. This makes generated questions more aligned
with pedagogical objectives and easier to evaluate.

A Bloom-aware question in the system requires at least these fields:

\begin{itemize}
    \item \texttt{question}: question content.
    \item \texttt{options}: list of answer options.
    \item \texttt{correct\_answer}: correct answer.
    \item \texttt{explanation}: explanation of why the answer is correct.
    \item \texttt{bloom\_level}: target Bloom level.
    \item \texttt{learning\_outcome\_id}: related Learning Outcome if available.
    \item \texttt{source\_chunk\_id}: source document chunk.
\end{itemize}

To improve quality, the system must validate LLM output after generation.
Basic checks include: the question has the required number of options,
there is exactly one correct answer, the explanation is grounded in the
source, the Bloom level belongs to the allowed value set, and the question
stem does not reveal the answer. Questions that do not satisfy these
checks can be moved to a review-needed state or sent back to the LLM for
schema-constrained repair.

\subsection{Micro-Content Classification in the System}
\label{subsec:phan-loai-micro-content}

In the proposed system, micro-content is classified by learning purpose
and by learner interaction pattern. This classification helps the content
generation pipeline produce stable outputs and helps the learning
interface present content in a clear rhythm.

\begin{longtable}{|p{3cm}|p{4.2cm}|p{6.8cm}|}
\caption{Micro-content classification in the system}
\label{tab:micro-content-types}\\
\hline
\textbf{Content type} & \textbf{Purpose} & \textbf{Characteristics} \\
\hline
\endfirsthead
\hline
\textbf{Content type} & \textbf{Purpose} & \textbf{Characteristics} \\
\hline
\endhead
Lesson Summary & Summarize a learning section & Presents the main ideas of a section or large chunk; usually longer than a card but shorter than a full lesson. \\
\hline
Key Concept Card & Clarify a central concept & Focuses on one concept, definition, example, or common misconception. \\
\hline
Flashcard & Support memory and review & The front side is a question/prompt and the back side is an answer or short explanation; suitable for spaced repetition. \\
\hline
Quiz Question & Check understanding & Question with answer, explanation, Bloom level, and source; used to evaluate learning progress. \\
\hline
Video Timestamp Card & Link content to video & Connects a card to start/end timestamps in a video or transcript. \\
\hline
\end{longtable}

All micro-content types must preserve source metadata to support
traceability, review, and regeneration when the original document changes.
This is an important difference between the proposed system and standalone
summarization or flashcard-generation tools.

\section{Large Language Models and Generative AI}
\label{sec:llm-generative-ai}

\subsection{Basic LLM Architecture}
\label{subsec:kien-truc-llm}

A Large Language Model (LLM) is a deep-learning model trained on large
amounts of text to learn language representations and generate token
sequences. The foundation architecture of most modern LLMs is the
Transformer, introduced by Vaswani et al.~\cite{vaswani2017attention}.
The core component of the Transformer is self-attention, which allows the
model to compute relevance between tokens in the input sequence and learn
long-range dependencies more effectively than traditional sequential
architectures.

LLMs usually operate through next-token prediction. Given an input
prompt, the model generates each new token based on the conditional
probability distribution of the previous context. In modern dialogue
models, the prompt is not only the user's question but also includes
system instructions, roles, retrieved context, output examples, and
format constraints.

In this thesis, LLMs are used for tasks that require knowledge
interpretation:

\begin{itemize}
    \item summarizing chunks into lesson summaries or key concept cards;
    \item generating flashcards and quiz questions;
    \item standardizing/extracting Learning Outcomes from course syllabi when needed;
    \item generating answers for the AI Tutor Chatbot based on RAG context;
    \item supporting answer explanations and the rationale for mapping questions to Bloom levels.
\end{itemize}

Models such as GPT-4o, Gemini, or open-weights models can be used
depending on cost, speed, privacy, and output-quality requirements
\cite{openai2024gpt4o,google2024gemini}. Therefore, the system design
places the LLM behind an abstract adapter or port so the provider can be
changed without changing business logic.

\subsection{Prompt Engineering}
\label{subsec:prompt-engineering}

Prompt Engineering is the technique of designing LLM inputs to control
behavior, quality, and output format without retraining the model. In an
educational system, a prompt should not only describe a generic task; it
must specify the pedagogical role, source material, learning objective,
Bloom constraint, output format, and safety conditions.

A quiz-generation prompt in the system usually contains:

\begin{enumerate}
    \item \textbf{Role}: ask the model to act as an instructor or
          pedagogical assistant.
    \item \textbf{Task}: generate multiple-choice questions, flashcards,
          or learning cards.
    \item \textbf{Source}: document chunk, heading path, Learning Outcome,
          and metadata.
    \item \textbf{Constraints}: number of questions, language, Bloom
          level, difficulty, maximum length.
    \item \textbf{Format}: JSON schema or fixed field structure.
    \item \textbf{Safety rules}: do not invent beyond the source, do not
          reveal the answer in the question, and return an error if
          context is insufficient.
\end{enumerate}

\begin{table}[H]
    \centering
    \caption{Prompt engineering techniques applied in the system}
    \label{tab:prompt-engineering-techniques}
    \begin{tabular}{p{3.4cm}p{4.6cm}p{5.5cm}}
        \toprule
        \textbf{Technique} & \textbf{Description} & \textbf{Application in this thesis} \\
        \midrule
        Role Prompting & Assigns a professional role to the model. & ``You are a university instructor preparing quizzes for learners''. \\
        Constraint Prompting & Sets constraints on quantity, length, difficulty, and language. & Generate exactly 5 questions, each with 4 options and a short explanation. \\
        Context Grounding & Provides retrieved context from documents. & Requires the answer to rely only on the supplied chunk. \\
        Few-shot Example & Provides an example of the desired output. & Shows a quiz item that satisfies schema and Bloom level. \\
        Self-check Instruction & Asks the model to check itself before responding. & Check unique correct answer, no answer leakage, and citation presence. \\
        \bottomrule
    \end{tabular}
\end{table}

Prompt Engineering improves output quality, but it is not sufficient to
guarantee correctness. The system still needs schema validation, source
data checks, access control, and instructor review.

\subsection{Structured Output Generation}
\label{subsec:structured-output-generation}

Structured Output Generation asks the LLM to return data in a defined
structure, for example JSON, instead of free-form text. In an LMS system,
structured output is important because AI results must be stored in a
database, displayed on the UI, edited by instructors, and evaluated by
automatic criteria.

If the LLM generates free-form text, the system has difficulty knowing
which part is the question, which part is the correct answer, which part
is the explanation, and where the source is. With a clear schema, the
backend can validate every field and reject invalid data. For example, a
quiz item can be modeled as:

\begin{lstlisting}[caption={Example output structure for a quiz item}]
{
  "type": "MCQ_SINGLE",
  "question": "Question ...",
  "options": ["A", "B", "C", "D"],
  "correct_answer": "B",
  "explanation": "Short explanation ...",
  "bloom_level": "Understand",
  "source_chunk_ids": ["chunk-123"]
}
\end{lstlisting}

After receiving the result, the system checks:

\begin{itemize}
    \item whether the JSON can be parsed;
    \item whether \texttt{type} is one of \texttt{MCQ\_SINGLE},
          \texttt{MCQ\_MULTI}, \texttt{TRUE\_FALSE},
          \texttt{FILL\_BLANK};
    \item whether the number and shape of \texttt{options} match
          \texttt{type} (for example, MCQ\_SINGLE needs $\geq 2$
          options);
    \item whether \texttt{correct\_answer} is a string for single-answer
          types and an array for \texttt{MCQ\_MULTI}, and references
          values that appear in \texttt{options};
    \item whether \texttt{bloom\_level} belongs to the allowed values;
    \item whether every id in \texttt{source\_chunk\_ids} exists in the
          course documents;
    \item whether the question, answer, and explanation are non-empty.
\end{itemize}

When data is invalid, the system can retry with a repair prompt or move
the result into a review-needed state. This approach reduces operational
errors caused by unstable LLM output.

\subsection{Hallucination in Educational Systems}
\label{subsec:hallucination-giao-duc}

Hallucination is the phenomenon in which an LLM generates information
that is incorrect, absent from the source, or inconsistent with context,
while still presenting it fluently and convincingly. In education,
hallucination is a major risk because learners may absorb incorrect
knowledge and instructors may spend time detecting and correcting errors.

Common causes include:

\begin{itemize}
    \item prompts that lack concrete source material;
    \item models reasoning from general knowledge instead of course documents;
    \item incorrect retrieval or missing important context;
    \item requests that ask for content beyond the available data;
    \item outputs that are not validated and reviewed before publishing.
\end{itemize}

In the proposed system, hallucination is reduced by these principles:

\begin{enumerate}
    \item \textbf{Retrieval-first}: every answer or generated content
          should be grounded in chunks retrieved from course materials.
    \item \textbf{Citation}: micro-content, quizzes, and chatbot answers
          must link sources to the original document/chunk.
    \item \textbf{Structured output}: output is validated before storage.
    \item \textbf{Human-in-the-loop}: instructors review before content is
          published to learners.
    \item \textbf{Refusal rule}: if no suitable context is found, the
          chatbot must say that the content is not present in the course
          instead of inventing an answer.
\end{enumerate}

\section{Retrieval-Augmented Generation (RAG)}
\label{sec:rag}

\subsection{RAG Architecture}
\label{subsec:kien-truc-rag}

Retrieval-Augmented Generation (RAG) combines information retrieval with
language generation \cite{lewis2020rag}. Instead of letting the LLM answer
entirely from knowledge learned during pretraining, the system retrieves
relevant document passages and injects them into the prompt as context.

A basic RAG architecture includes five steps:

\begin{enumerate}
    \item \textbf{Ingestion}: documents are parsed, chunked, embedded, and
          stored in a vector database.
    \item \textbf{Query understanding}: the user's question or content
          generation request is normalized.
    \item \textbf{Retrieval}: the system finds relevant top-K chunks in
          the vector store, possibly filtering by course, document, role,
          or Learning Outcome.
    \item \textbf{Prompt assembly}: relevant chunks are inserted into the
          prompt along with instructions and output schema.
    \item \textbf{Generation}: the LLM generates answers, micro-content,
          or quizzes based on retrieved context.
\end{enumerate}

\begin{figure}[H]
    \centering
    \begin{tikzpicture}[
        node distance=1.2cm,
        box/.style={draw, rounded corners, align=center, minimum width=3.1cm, minimum height=1cm},
        arrow/.style={->, thick}
    ]
        \node[box] (doc) {Documents\\PDF/DOCX/PPTX};
        \node[box, right=of doc] (chunk) {Parse\\Chunk\\Embedding};
        \node[box, right=of chunk] (store) {Vector DB\\Metadata};
        \node[box, below=of store] (retrieve) {Retrieve\\top-K chunks};
        \node[box, left=of retrieve] (prompt) {Prompt\\+ Context};
        \node[box, left=of prompt] (llm) {LLM\\Generation};
        \draw[arrow] (doc) -- (chunk);
        \draw[arrow] (chunk) -- (store);
        \draw[arrow] (store) -- (retrieve);
        \draw[arrow] (retrieve) -- (prompt);
        \draw[arrow] (prompt) -- (llm);
    \end{tikzpicture}
    \caption{General RAG architecture in the system}
    \label{fig:rag-architecture}
\end{figure}

In educational systems, RAG is especially important because course
materials are the source of truth. For the same question, the appropriate
answer must follow the course content, the instructor's definitions, and
the curriculum scope.

\subsection{Semantic Retrieval}
\label{subsec:semantic-retrieval}

Semantic Retrieval retrieves documents by meaning rather than exact
keyword matching. Each chunk and each query are represented as embedding
vectors in a high-dimensional space. Texts with similar meaning should
have nearby vectors according to a metric such as cosine similarity.

With an embedding model \(f\), a chunk \(c_i\) is represented as:

\[
    \vec{e_i} = f(c_i)
\]

A query \(q\) is represented as:

\[
    \vec{e_q} = f(q)
\]

Cosine similarity between query and chunk is calculated as:

\[
    sim(\vec{e_q}, \vec{e_i}) =
    \frac{\vec{e_q} \cdot \vec{e_i}}{\|\vec{e_q}\| \|\vec{e_i}\|}
\]

The system selects chunks with the highest similarity for context. To
avoid retrieving the wrong scope, semantic retrieval must be combined
with metadata filtering: \texttt{course\_id}, \texttt{document\_id},
publish state, access permissions, and Learning Outcome if available.
Multilingual embedding models such as BGE-M3 can support Vietnamese and
English retrieval in the same representation space \cite{chen2024bge}.

\subsection{Curriculum-aware Retrieval}
\label{subsec:curriculum-aware-retrieval}

Curriculum-aware Retrieval extends semantic retrieval for educational
contexts. Instead of only finding the chunks closest to a query, the
system also considers curriculum structure such as Course, Chapter,
Learning Outcome, Bloom level, and Assessment. This helps retrieve the
right knowledge that instructors want to assess or learners need to
review.

For example, if an instructor requests quiz generation for LO-3.2, the
system should not take every chunk with similar keywords. It should
prioritize chunks mapped to LO-3.2, located in the relevant chapter, and
already reviewed. If a learner is weak in a specific LO, the system can
find cards, flashcards, and quizzes related to that LO for review.

A curriculum-aware pipeline can include:

\begin{enumerate}
    \item Extract or import the list of Learning Outcomes from the course
          syllabus.
    \item Map document chunks to Learning Outcomes using headings,
          keywords, embedding similarity, or LLMs.
    \item Store chunk--LO relationships in the metadata store or graph
          database.
    \item During retrieval, filter or boost chunks belonging to the target
          LO.
    \item Generate content/quiz and record LO, Bloom level, and source
          chunks for analytics.
\end{enumerate}

Curriculum-aware Retrieval is the foundation for Bloom-aware Quiz
Generation, weak-LO dashboards, and adaptive learning in later phases.

\subsection{GraphRAG}
\label{subsec:graphrag}

GraphRAG combines RAG with a Knowledge Graph. In traditional RAG, the main
retrieval unit is a text chunk. In GraphRAG, the system adds nodes and
edges representing relationships among Course, Chapter, Learning Outcome,
Concept, Chunk, Quiz, and Assessment. Retrieval can then rely not only on
similarity but also on structured relationships.

In the proposed system, GraphRAG is understood in a curriculum-aware
sense; it does not require extracting all world knowledge into a graph.
Important relationships include:

\begin{itemize}
    \item Course contains Chapter.
    \item Chapter has Learning Outcome.
    \item Chunk supports Learning Outcome.
    \item Quiz measures Learning Outcome.
    \item Concept appears in Chunk.
\end{itemize}

\begin{table}[H]
    \centering
    \caption{Comparison between traditional RAG and GraphRAG in educational systems}
    \label{tab:rag-vs-graphrag-c3}
    \begin{tabular}{p{3.8cm}p{4.8cm}p{5.2cm}}
        \toprule
        \textbf{Criterion} & \textbf{Traditional RAG} & \textbf{GraphRAG / Curriculum-aware RAG} \\
        \midrule
        Retrieval unit & Text chunk & Chunk combined with nodes/edges such as LO, Chapter, Concept \\
        Search method & Vector similarity & Vector similarity + metadata filter + graph traversal \\
        Alignment with outcomes & Depends on simple metadata & Explicit relation between Chunk and Learning Outcome \\
        Quiz by LO & May drift if query is ambiguous & Prioritizes chunks supporting the target LO \\
        Implementation complexity & Lower & Higher because graph schema and data synchronization are needed \\
        \bottomrule
    \end{tabular}
\end{table}

GraphRAG does not replace semantic retrieval; it adds a structural layer.
In the MVP, the system can use metadata and chunk--LO relationships at a
basic level. In future expansion, the graph can rerank results, explain
why a chunk was selected, or generate quizzes that connect multiple
concepts.

\section{Text Chunking and Document Intelligence}
\label{sec:text-chunking-document-intelligence}

\subsection{Fixed-size Chunking}
\label{subsec:fixed-size-chunking}

Fixed-size Chunking divides text by fixed length, for example 500 or 1000
tokens per chunk. This is simple to implement, easy to control in cost,
and useful when documents have no clear structure.

Its advantage is that chunk sizes are relatively uniform, which helps
predict embedding cost and the number of tokens sent to the LLM. Its main
drawback is that it can cut across sentences, tables, or complete ideas.
A concept may be split into two chunks, reducing retrieval quality and
leaving the LLM without enough context when generating content.

In the proposed system, fixed-size chunking is suitable as a fallback when
documents have no headings, the parser cannot extract structure, or the
input is plain text.

\subsection{Heading-based Chunking}
\label{subsec:heading-based-chunking}

Heading-based Chunking divides documents by structural headings such as
chapter, section, subsection, or slide title. This is suitable for
academic materials because content is usually organized by topic. When
chunks follow headings, each chunk is more likely to contain a relatively
complete knowledge unit.

For example, a document with this structure:

\begin{itemize}
    \item Chapter 1: Overview
    \item 1.1 Basic concepts
    \item 1.2 Illustrated examples
\end{itemize}

The system can create chunks by section and store \texttt{heading\_path}
to know the chunk's location. This metadata is important for citation,
review UI, and chapter-based retrieval.

The limitation of heading-based chunking is its dependence on source
document quality. If a PDF has no clear heading structure, or slides use
inconsistent styles, the system must fall back to paragraph-, sentence-,
or size-based chunking.

\subsection{Semantic Chunking}
\label{subsec:semantic-chunking}

Semantic Chunking divides text based on changes in meaning between
paragraphs. Instead of relying only on token count or headings, the
system analyzes sentence/paragraph embeddings to find split points when
the topic changes. In theory, semantic chunking creates more natural
chunks and preserves semantic integrity better.

However, semantic chunking is more expensive because it requires
sentence- or paragraph-level embeddings before deciding boundaries. It is
also harder to debug: instructors and developers cannot easily predict
why the system split at a specific location. Therefore, in this system,
semantic chunking is more suitable as a later quality optimization, while
the MVP prioritizes heading-based chunking with fallback.

\subsection{Sliding Window and Overlap}
\label{subsec:sliding-window-overlap}

Sliding Window divides text into fixed-size windows that move by a fixed
step. Overlap is the repeated part between two consecutive chunks. For
example, if chunk size is 800 tokens and overlap is 100 tokens, the next
chunk repeats the final 100 tokens of the previous chunk.

Overlap reduces the risk of losing context at chunk boundaries. This is
important in academic documents, because a definition may appear at the
end of one section while its example appears at the beginning of the next.
Without overlap, retrieval may miss relevant information.

However, overlap also increases token count, embedding cost, and duplicate
retrieval results. Therefore, overlap must be just enough. In the proposed
system, sliding window is mainly used when a heading section is too long
or when falling back for unstructured documents.

\subsection{Trade-off in Chunk Design}
\label{subsec:tradeoff-thiet-ke-chunk}

Chunk design is one of the most important decisions in a RAG pipeline. If
chunks are too short, retrieval can be precise at sentence level but may
lack enough context for the LLM to generate answers or quizzes. If chunks
are too long, they preserve more context but can dilute content, increase
token cost, and make semantic search less sharp.

\begin{longtable}{|p{3cm}|p{4.2cm}|p{6.8cm}|}
\caption{Trade-offs in chunk design}
\label{tab:chunking-tradeoff}\\
\hline
\textbf{Factor} & \textbf{If too low/small} & \textbf{If too high/large} \\
\hline
\endfirsthead
\hline
\textbf{Factor} & \textbf{If too low/small} & \textbf{If too high/large} \\
\hline
\endhead
Chunk size & Lacks context, hard to generate questions with sufficient explanation. & Dilutes content, reduces retrieval precision, increases LLM tokens. \\
\hline
Overlap & Information can be lost at chunk boundaries. & More chunks, higher embedding cost, duplicate results. \\
\hline
Top-K count & Relevant chunks may be missed. & Long prompt, noisy context, higher cost and latency. \\
\hline
Heading depth & Chunks are too broad if only high-level headings are used. & Chunks are too fragmented if split by too small headings. \\
\hline
Source metadata & Citation and review are difficult if metadata is poor. & Overly complex metadata increases storage and synchronization cost. \\
\hline
\end{longtable}

For academic documents, a suitable strategy is to prioritize
heading-based chunking and then apply sliding window for overly long
sections. Each chunk must store at least \texttt{document\_id},
\texttt{chunk\_id}, \texttt{heading\_path}, page or timestamp position,
language, and publish/review state. These metadata are the foundation for
RAG, GraphRAG, curriculum-aware retrieval, and human-in-the-loop review.

\section{Vector Embedding and Semantic Search}
\label{sec:vector-embedding-semantic-search}

Vector embedding represents text as numeric vectors in a high-dimensional
space. Instead of direct keyword matching, the system compares the meaning
of a user question and document chunks through vector distance or
similarity. This approach is the foundation of semantic search, dense
retrieval, and many modern RAG architectures
\cite{karpukhin2020dpr,lewis2020rag}.

In the proposed LMS system, embeddings are generated for document chunks,
video transcripts, published flashcards/quizzes, and user queries. When a
learner searches or asks the AI Tutor Chatbot, the system converts the
query into a vector, retrieves the most relevant chunks, and uses them as
the basis for displaying results or providing context to the LLM.
Therefore, embedding quality directly affects semantic search, citation,
chatbot answers, and source-grounded quiz generation.

\subsection{Embedding Models}
\label{subsec:embedding-models}

An embedding model receives text and maps it to a fixed-dimensional vector.
Modern semantic representation models are often based on Transformers,
inherit contextual-representation ideas from BERT, and are further
optimized for sentence comparison, passage retrieval, or document search
\cite{devlin2019bert,reimers2019sbert}.

For semantic search, embeddings should place semantically similar passages
near each other even if they do not share the same keywords. For example,
the query ``how the system generates questions based on output standards''
should retrieve chunks about Learning Outcome, Bloom level, and quiz
generation even if the document does not contain that exact phrase.

Two model groups are common. The first is bi-encoder, where query and
document are encoded independently; this is suitable for large-scale
search because document vectors can be computed ahead of time and stored
in a vector database. The second is cross-encoder, where query and
document are passed into the model together for relevance scoring; it is
usually more accurate but slower, so it is better for reranking a small
top-K result set instead of searching the whole corpus.

In the proposed system, bi-encoder is the main choice for the MVP because
it fits the asynchronous pipeline: when a document or transcript is
chunked, the system creates embeddings once and stores them in Qdrant.
Cross-encoder or LLM reranker can be added later to improve results for
chatbot or Learning Outcome-based quiz generation.

\subsection{Similarity Search}
\label{subsec:similarity-search}

Similarity search finds vectors closest to the query vector. For each
query, the system creates an embedding, compares it with stored chunk
vectors, and returns the top-K results with the highest similarity. Common
metrics include cosine similarity, dot product, and Euclidean distance.
The metric choice depends on how the embedding model was trained and how
vectors are normalized.

Cosine similarity measures the angle between two vectors and is useful
when comparing semantic direction rather than vector magnitude. Dot
product is often used when vectors are normalized or the model is trained
directly for inner-product scoring. Euclidean distance measures geometric
distance between two points in vector space. Although the formulas differ,
the common goal is to rank chunks by relevance to the query.

In an LMS, similarity search should not rely only on vector scores.
Results must combine metadata filters such as course, chapter, publish
state, user role, and Learning Outcome. For example, learners can only
retrieve published content in their own course, while instructors may also
search draft content for review. Combining vector similarity and metadata
filtering makes semantic search both meaningful and compliant with
business rules.

\subsection{Approximate Nearest Neighbor}
\label{subsec:approximate-nearest-neighbor}

When the number of chunks is small, exact search can compare the query
vector with every vector in the database. However, when the number grows
to hundreds of thousands or millions, linear search becomes expensive and
cannot meet latency requirements. Approximate Nearest Neighbor (ANN)
solves this by finding approximate nearest vectors at lower cost than
exact search.

Common ANN approaches include inverted file indexes, product quantization,
and graph-based indexes such as HNSW. HNSW builds a multi-layer graph in
which search moves through nodes progressively closer to the query vector.
Many modern vector databases use or support HNSW because it balances query
speed, retrieval accuracy, and scalability
\cite{malkov2018hnsw,johnson2019faiss}.

ANN always has a trade-off between speed and accuracy. If the configuration
prioritizes speed, the system may miss some relevant chunks; if it
prioritizes accuracy, latency and compute cost increase. Therefore, in the
proposed system, top-K, index type, and search configuration should be
evaluated with a test query set instead of chosen by intuition. For the
MVP, Qdrant acts as the vector database storing embeddings, metadata, and
ANN retrieval for documents and video transcripts.

\subsection{Retrieval Metrics}
\label{subsec:retrieval-metrics}

Retrieval metrics evaluate whether semantic search retrieves the correct
sources. This is important because if retrieval is wrong, the LLM in RAG
may answer without evidence or generate quizzes from unsuitable chunks.
Benchmarks such as BEIR and MTEB show that retrieval evaluation should
cover many tasks, data domains, and query types
\cite{thakur2021beir,muennighoff2023mteb}.

Common metrics include Recall@K, Precision@K, MRR, and nDCG. Recall@K
measures whether the correct document appears in the first K results; it
is suitable for RAG because the LLM only needs to receive at least a few
correct chunks in context. Precision@K measures how many top-K results
are actually relevant, helping detect noisy context. MRR emphasizes the
rank of the first correct result, while nDCG considers both ranking and
multi-level relevance.

\begin{longtable}{|p{3cm}|p{4.2cm}|p{6.8cm}|}
\caption{Retrieval evaluation metrics in semantic search}
\label{tab:retrieval-metrics}\\
\hline
\textbf{Metric} & \textbf{Meaning} & \textbf{Application in the system} \\
\hline
\endfirsthead
\hline
\textbf{Metric} & \textbf{Meaning} & \textbf{Application in the system} \\
\hline
\endhead
Recall@K & Checks whether the correct document/chunk is in top-K. & Evaluates whether chatbot or quiz generator receives enough relevant sources. \\
\hline
Precision@K & Measures the proportion of truly relevant top-K results. & Reduces noisy context and limits wrong chunks in prompts. \\
\hline
MRR & Measures the rank of the first correct result. & Prioritizes the best source in search UI or RAG context. \\
\hline
nDCG@K & Evaluates ranking when results have multiple relevance levels. & Suitable when instructors grade relevance as high, medium, low. \\
\hline
\end{longtable}

For this thesis, the retrieval evaluation set should include sample
questions for each course, correct chunks labeled by instructors or
developers, embedding/index/top-K configurations, and measured Recall@K,
Precision@K, and MRR. These metrics evaluate semantic search itself and
also support comparisons when changing chunking strategy, embedding model,
reranker, or vector database configuration.

\section{Learning Management System Integration Standards}
\label{sec:lms-integration-standards}

In the digital-education ecosystem, a learning tool rarely works alone.
It usually needs to integrate with an LMS such as Moodle, Canvas, or Open
edX to receive course, user, role, assignment, and grade information. If
each LMS required a separate integration mechanism, development and
maintenance cost would increase significantly. Therefore, integration
standards such as Learning Tools Interoperability (LTI), OAuth2, and
OpenID Connect are used to standardize how LMS platforms connect with
external learning tools \cite{lti-advantage,openedx-lti}.

For this thesis, the micro-content, quiz, and AI Tutor Chatbot system is
treated as an external learning tool. The LMS is the main platform for
managing courses and learners, while the proposed system provides
specialized AI functions. This approach prevents dependence on a single
LMS while still allowing the system to receive course context, enforce
user authorization, and return learning results to the LMS when needed.

\subsection{Learning Tools Interoperability (LTI)}
\label{subsec:lti}

Learning Tools Interoperability (LTI) is a standard developed by 1EdTech
to allow LMS platforms to connect with external learning tools. In the
LTI model, the LMS is usually called the platform, and the external
application is called the tool. When a user clicks a link inside the LMS,
the platform initiates a launch and sends the tool the necessary
information such as user identity, course, role, and resource link.

LTI solves three main problems. First, it provides single sign-on so users
do not need to log in again to the tool. Second, it transmits learning
context so the tool knows which course, lesson, or activity the user is
currently in. Third, it standardizes LMS--tool communication, allowing the
same tool to integrate with multiple LMS platforms if they support LTI.

In the proposed system, LTI launch opens the learner interface or the
instructor interface from the LMS. After a valid launch, the system
creates an internal session, maps the LMS user to a system account, and
determines access rights. For example, instructors can upload documents,
review, and publish content; learners can only view published content,
take quizzes, and interact with the chatbot within the course scope.

\subsection{LTI 1.3 and Advantage Services}
\label{subsec:lti13-advantage-services}

LTI 1.3 is the modern version of LTI, using security mechanisms based on
OAuth2, OpenID Connect, and JSON Web Token. Compared with older LTI
versions based on OAuth 1.0, LTI 1.3 better standardizes login, signing,
claims in ID tokens, and registration information between platform and
tool.

LTI Advantage is a set of extended services built on LTI 1.3. These
services allow a tool not only to be launched from the LMS but also to
interact more deeply with it. Three important capability groups are Deep
Linking, which allows the tool to create or select content to place in
the LMS; Assignment and Grade Services, which synchronize assignments,
scores, and submission states; and Names and Role Provisioning Services,
which retrieve membership lists and roles in a learning context.

\begin{longtable}{|p{3.5cm}|p{4.2cm}|p{6.3cm}|}
\caption{LTI Advantage services related to the system}
\label{tab:lti-advantage-services}\\
\hline
\textbf{Service} & \textbf{Purpose} & \textbf{Application in this thesis} \\
\hline
\endfirsthead
\hline
\textbf{Service} & \textbf{Purpose} & \textbf{Application in this thesis} \\
\hline
\endhead
Deep Linking & Select or create content from the tool to attach to the LMS. & Instructors insert flashcard sets, quizzes, or chatbot activities into a course. \\
\hline
Assignment and Grade Services & Synchronize assignments, grades, and result states. & Send quiz scores or micro-learning completion states back to the LMS. \\
\hline
Names and Role Provisioning Services & Retrieve learners and roles in a context. & Support authorization, class analytics, and learner-list synchronization. \\
\hline
\end{longtable}

In the MVP, the system can prioritize LTI launch and Deep Linking for the
basic learning flow. AGS and NRPS can be implemented incrementally when
grade synchronization, class rosters, or learning dashboards need LMS
integration.

\subsection{OAuth2 and OpenID Connect in LMS}
\label{subsec:oauth2-oidc-lms}

OAuth2 is an authorization framework that allows an application to access
resources on behalf of a user or system within granted scopes. OpenID
Connect (OIDC) extends OAuth2 to support identity authentication, usually
returning the authentication result as an ID token. LTI 1.3 uses OIDC to
perform the login flow between LMS and tool, and uses JWT to transmit
signed claims.

During LTI launch, the tool should not trust data sent from the browser
before verification. The tool must validate issuer, audience, nonce,
message type, deployment id, JWT signature, and token expiration. These
checks ensure that the launch request really comes from the registered
LMS, has not been replayed, and is intended for the correct tool.

In the proposed system, after OIDC verification succeeds, the backend
creates a server-side internal session. This session stores necessary
information such as user id, role, course id, resource link id, and
expiration. Business APIs such as upload, review, publish, search,
chatbot, and quiz submission rely on this session for access control
instead of trusting client-submitted data.

\subsection{Deep Linking}
\label{subsec:deep-linking}

Deep Linking is an LTI Advantage service that allows instructors to
select or create content in the tool and return it to the LMS as a
learning link. Instead of manually copying URLs or configuring resources,
the LMS can open the tool in content-selection mode; the tool then returns
one or more content items for the LMS to attach to the course.

In the proposed system, Deep Linking fits the instructor AI-material
creation flow. After uploading documents and reviewing content, the
instructor can select a flashcard set, quiz, lesson card set, or chatbot
activity to insert into an LMS module. When learners open that activity,
the LMS performs an LTI launch with the corresponding resource link, so
the system knows what content to display.

Deep Linking also reduces UI dependence. The LMS remains the main place
for course organization and navigation, while the AI system is responsible
for creating, managing, and distributing specialized content. This aligns
with the thesis goal of building an AI tool that can be integrated into
an LMS instead of replacing the entire LMS.

\subsection{Assignment and Grade Services (AGS)}
\label{subsec:assignment-grade-services}

Assignment and Grade Services (AGS) allow a tool to exchange score and
submission-state information with the LMS. On the LMS side, an activity
can have a corresponding line item in the gradebook. On the tool side,
after a learner completes a quiz or learning activity, the tool can send
the score back to the LMS through that line item.

In this thesis, AGS can be used for automatically generated quizzes or
micro-learning completion activities. For example, when a learner takes a
quiz, the system stores an internal submission containing answers, score,
time, and question sources. If the activity is linked to the LMS
gradebook, the backend sends the score back through AGS. Instructors can
therefore continue viewing grades in the familiar LMS, while the AI system
retains richer data for analytics.

It is important to distinguish grading data from learning-analytics data.
The LMS usually only needs final score, completion state, or submission
time. The proposed system needs to store additional details such as
questions, Bloom levels, Learning Outcomes, attempt count, response time,
and answer explanations to evaluate learning quality. Therefore, AGS is
the official grade-synchronization channel, not a replacement for the
internal analytics database.

\subsection{LMS Interoperability Architecture}
\label{subsec:lms-interoperability-architecture}

LMS integration architecture must clearly separate LMS communication
standards from core business logic. If upload, review, semantic search, or
quiz logic depends directly on one LMS-specific API, the system becomes
difficult to extend to Moodle, Canvas, Open edX, or other platforms.
Therefore, the system should place LTI/OIDC/AGS in an adapter layer, while
the domain/application core works with abstract concepts such as user,
role, course, activity, and grade result.

A suitable integration architecture for this thesis includes:

\begin{itemize}
    \item \textbf{LMS Platform:} Moodle, Canvas, or Open edX, where
          courses, users, and learning navigation are managed.
    \item \textbf{LTI/OIDC Adapter:} handles login initiation, launch
          validation, JWT claims, deployment configuration, and session
          mapping.
    \item \textbf{Deep Linking Adapter:} creates content items so the LMS
          can insert flashcards, quizzes, or chatbot activities into a
          course.
    \item \textbf{AGS Adapter:} synchronizes quiz score, completion
          state, or score back to the LMS gradebook.
    \item \textbf{Application Core:} contains use cases for upload,
          processing, review, publish, semantic search, chatbot, quiz, and
          analytics.
\end{itemize}

\begin{table}[H]
    \centering
    \caption{Data mapping between LMS and the proposed system}
    \label{tab:lms-data-mapping}
    \begin{tabular}{p{4cm}p{4.2cm}p{5.3cm}}
        \toprule
        \textbf{Data from LMS} & \textbf{Object in system} & \textbf{Usage purpose} \\
        \midrule
        User claim & User / Account & Identify identity and create session. \\
        Role claim & Role / Permission & Authorize instructors, learners, and administrators. \\
        Context claim & Course / Class context & Limit content by course. \\
        Resource link & Learning activity & Open the correct flashcard, quiz, or chatbot activity. \\
        Line item & Grade mapping & Send score or completion state through AGS. \\
        \bottomrule
    \end{tabular}
\end{table}

With this organization, LMS integration becomes a boundary layer of the
system. When changing LMS or adding a new platform, the team mainly
updates adapters and deployment configuration, while core AI use cases
remain stable. This is an important condition for deploying the system in
many training environments without rewriting the document processing,
semantic search, or AI Tutor Chatbot pipelines.

\section{Knowledge Graph and Curriculum Mapping}
\label{sec:knowledge-graph-curriculum-mapping}

In an AI-based learning system, content does not only exist as isolated
text chunks. Each document passage, video transcript, flashcard, or quiz
question is usually related to a chapter, a concept, a Learning Outcome,
and a specific assessment activity. Knowledge Graphs model these entities
and their relationships as nodes and edges, supporting retrieval, source
explanation, and structured learning-progress analytics
\cite{hogan2021knowledgegraphs}.

Curriculum Mapping links learning content with curriculum structure,
especially Learning Outcomes, topics, and assessments. In this thesis,
Curriculum Mapping helps the system understand which Outcome a chunk
supports, which competency a quiz measures, and which knowledge group a
learner is weak in. It is the foundation for Bloom-level quiz generation,
learning analytics dashboards, and later adaptive learning. This approach
aligns with constructive alignment and curriculum mapping in educational
design, where learning objectives, learning activities, and assessment
must be explicitly connected
\cite{biggs1996constructive,harden2001curriculum}.

\subsection{Learning Outcome Graph}
\label{subsec:learning-outcome-graph}

The Learning Outcome Graph represents relationships among course,
chapter, learning outcomes, concepts, document chunks, and assessment
questions. Instead of storing a Learning Outcome as an isolated text
field, the system models it as a node connected to other nodes. This
representation fits education because learning content is hierarchical
and contains many cross-relations.

In the proposed system, important nodes include \texttt{Course},
\texttt{Chapter}, \texttt{LearningOutcome}, \texttt{Concept},
\texttt{Chunk}, \texttt{Flashcard}, \texttt{QuizQuestion}, and
\texttt{Assessment}. Edges express relations such as course contains
chapter, chapter has Learning Outcome, chunk supports Learning Outcome,
quiz measures Learning Outcome, or concept appears in chunk. These
relations can be created from the course syllabus, document headings,
user-provided metadata, and automatic extraction results.

\begin{longtable}{|p{3.2cm}|p{4.2cm}|p{6.6cm}|}
\caption{Main relationships in the Learning Outcome Graph}
\label{tab:learning-outcome-graph-relations}\\
\hline
\textbf{Relationship} & \textbf{Meaning} & \textbf{Application in the system} \\
\hline
\endfirsthead
\hline
\textbf{Relationship} & \textbf{Meaning} & \textbf{Application in the system} \\
\hline
\endhead
Course--Chapter & A course consists of many chapters or modules. & Filter content by LMS course structure. \\
\hline
Chapter--LearningOutcome & Each chapter supports one or more outcomes. & Generate quizzes and report results by LO. \\
\hline
Chunk--LearningOutcome & A document chunk provides evidence for an LO. & Retrieve the right source when chatbot answers or quizzes are generated. \\
\hline
QuizQuestion--LearningOutcome & A question assesses one or more LOs. & Analyze learner weaknesses by outcome. \\
\hline
Concept--Chunk & A concept appears in one or more chunks. & Support search, review recommendation, and GraphRAG expansion. \\
\hline
\end{longtable}

The Learning Outcome Graph does not have to be complete in the MVP. The
first phase can start with basic relations among course, chapter, LO, and
chunk. As the system collects quiz submissions, chatbot interactions, and
instructor feedback, the graph can be enriched to reflect the knowledge
structure and learner mastery more accurately.

\subsection{Curriculum-aware Learning}
\label{subsec:curriculum-aware-learning}

Curriculum-aware Learning organizes the learning experience according to
curriculum structure and Learning Outcomes, instead of only relying on a
document list or upload order. When the system knows which content belongs
to which LO, it can help learners study toward specific objectives and
help instructors evaluate outcome achievement.

In the proposed system, Curriculum-aware Learning appears in three main
flows. First, when generating flashcards and quizzes, the prompt can
include LO, Bloom level, and chapter so the LLM creates pedagogically
aligned content. Second, when learners search or ask the chatbot,
retrieval can prioritize chunks belonging to the LO or chapter being
studied. Third, learning analytics can aggregate quiz results by LO to
identify knowledge groups that the learner or class needs to review.

A typical example is quiz generation by outcome. The instructor selects
an LO in the syllabus, the system retrieves chunks supporting that LO, and
then asks the LLM to generate questions at the appropriate Bloom level.
Each question is stored with \texttt{lo\_id}, \texttt{bloom\_level},
\texttt{source\_chunk\_id}, and review state. When learners take the quiz,
scores can be aggregated by LO to build a progress dashboard.

The curriculum-aware approach also reduces the risk of generating content
that drifts away from the objective. If the system relies only on general
similarity search, a short query may retrieve chunks about a similar topic
but not the correct outcome. Combining curriculum metadata allows the
system to limit or boost sources that better match teaching objectives.

\subsection{Metadata Strategy}
\label{subsec:metadata-strategy}

Metadata Strategy defines what information the system must store for
traceability, retrieval, review, publishing, and AI-quality evaluation.
In RAG and curriculum mapping, metadata is as important as text content
because it records where a chunk comes from, which course it belongs to,
which LO it relates to, and whether learners are allowed to see it.

In the proposed system, metadata should be organized into multiple layers.
The source layer includes \texttt{document\_id}, \texttt{file\_type},
\texttt{page}, \texttt{timestamp}, or \texttt{heading\_path}. The
pedagogical layer includes \texttt{course\_id}, \texttt{chapter\_id},
\texttt{learning\_outcome\_id}, \texttt{bloom\_level}, and
\texttt{difficulty}. The operational layer includes
\texttt{pipeline\_status}, \texttt{parser\_version},
\texttt{embedding\_model}, \texttt{prompt\_version},
\texttt{review\_status}, and \texttt{publish\_status}.

\begin{table}[H]
    \centering
    \caption{Metadata strategy for curriculum-aware retrieval}
    \label{tab:metadata-strategy-curriculum}
    \footnotesize
    \begin{tabularx}{\textwidth}{
        >{\raggedright\arraybackslash}p{2.7cm}
        >{\raggedright\arraybackslash}p{4.2cm}
        >{\raggedright\arraybackslash}X}
        \toprule
        \textbf{Metadata group} & \textbf{Representative fields} & \textbf{How the system uses it} \\
        \midrule
        Source location
            & \texttt{source\_type}, \texttt{document\_id}, \texttt{page},
              \texttt{heading\_path}, \texttt{start\_ms}, \texttt{end\_ms}
            & Shows citations in generated cards, quiz explanations, and
              chatbot answers; lets instructors open the original page or
              video timestamp during review. \\
        \midrule
        Curriculum scope
            & \texttt{course\_id}, \texttt{chapter\_id}, \texttt{lo\_id},
              \texttt{assessment\_id}
            & Applies hard filters so retrieval stays inside the launched
              Canvas course and the selected chapter/LO; prevents chunks from
              unrelated courses from entering the RAG context. \\
        \midrule
        Pedagogical intent
            & \texttt{bloom\_level}, \texttt{difficulty},
              \texttt{content\_type}, \texttt{learning\_goal}
            & Guides quiz generation, distractor design, and micro-content
              type selection; supports later analysis of learner performance
              by LO and Bloom level. \\
        \midrule
        AI operation
            & \texttt{parser\_version}, \texttt{chunker\_version},
              \texttt{embedding\_model}, \texttt{prompt\_version}
            & Records which pipeline version produced a chunk or generated
              item, making evaluation reproducible after changing models,
              prompts, or chunking strategy. \\
        \midrule
        Access and lifecycle
            & \texttt{review\_status}, \texttt{publish\_status},
              \texttt{visibility}, \texttt{updated\_at}
            & Controls whether a chunk can be shown to learners, included in
              chatbot retrieval, or restricted to instructor-only review
              before publication. \\
        \bottomrule
    \end{tabularx}
\end{table}

Metadata must be stored consistently in both the metadata store and the
vector database. The vector database uses metadata for retrieval filters,
while PostgreSQL or graph store keeps more detailed relationships for
review, analytics, and audit. The key principle is that every AI-generated
content item must be traceable back to its source: document, chunk,
prompt, model, and pipeline version.

\subsection{Graph-based Retrieval}
\label{subsec:graph-based-retrieval}

Graph-based Retrieval combines vector similarity with graph relationships.
Vector search finds chunks semantically close to the query, while the
graph expands or filters results based on relationships such as same
Learning Outcome, same chapter, same concept, or same assessment. This is
suitable for educational systems because learner questions often depend
on a specific course context, not only on general linguistic meaning.

A graph-based retrieval flow can have four steps. First, the system
identifies context from LMS or UI, for example course, chapter, resource
link, or target LO. Next, the query is embedded and top-K chunks are
retrieved from the vector database. Then, the system uses the graph to
filter, expand, or rerank chunks according to curriculum relationships.
Finally, chunks are sent to the chatbot or quiz generator prompt with
citations and source metadata.

Graph-based retrieval can be implemented at several levels. At the simple
level, the system uses metadata filters such as \texttt{course\_id},
\texttt{chapter\_id}, and \texttt{publish\_status}. At a more advanced
level, it queries the graph for chunks belonging to related LOs, nearby
concepts, or questions that previously assessed the same outcome. At a
complete level, the graph supports reranking and explanation of why a
source was selected.

\begin{longtable}{|p{3.2cm}|p{4.5cm}|p{6.3cm}|}
\caption{Role of graph in retrieval flows}
\label{tab:graph-based-retrieval-roles}\\
\hline
\textbf{Flow} & \textbf{Graph usage} & \textbf{Benefit} \\
\hline
\endfirsthead
\hline
\textbf{Flow} & \textbf{Graph usage} & \textbf{Benefit} \\
\hline
\endhead
Semantic Search & Filter by course, chapter, LO, and publish state. & Results match learning context and access permissions. \\
\hline
AI Tutor Chatbot & Expand context with chunks from the same LO or related concepts. & Answers become more complete while remaining verifiable. \\
\hline
Quiz Generation & Prioritize chunks supporting the target LO and desired Bloom level. & Questions align with outcomes and are easier to review/evaluate. \\
\hline
Learning Analytics & Aggregate submissions by LO, chapter, and concept. & Detect difficult content and recommend evidence-based review. \\
\hline
\end{longtable}

In the MVP, graph-based retrieval can begin with metadata filters and
basic chunk--LO relations. As data and evaluation needs increase, the
system can add a graph database or graph index to support more complex
relationship queries. The important point is to design the data from the
beginning to preserve \texttt{source\_chunk\_id},
\texttt{learning\_outcome\_id}, and related metadata, so that future
GraphRAG expansion does not require reprocessing all data.

\section{Speech-to-Text and Video AI}
In a modern LMS enhanced by Generative AI, Speech-to-Text and Video AI
are foundational technologies. They digitize, extract, and structure
multimedia data (lecture video and audio) into text streams and metadata
that Large Language Models can understand and process, thereby enabling
micro-content and quiz generation.

\subsection{Whisper}
Whisper is an advanced open-source Automatic Speech Recognition (ASR)
system based on the Transformer architecture developed by OpenAI
\cite{radford2023robust}.
\begin{itemize}
    \item \textbf{Mechanism:} The model is trained on hundreds of
          thousands of hours of multilingual audio, allowing it to
          transcribe audio into text with high accuracy. It handles
          background noise, diverse accents, and complex technical terms
          well.
    \item \textbf{Role in LMS:} Whisper digitizes the entire spoken
          lecture into raw transcript and provides sentence-level
          timestamps, creating the foundation for later semantic
          processing.
\end{itemize}

\subsection{Transcript Alignment}
Transcript Alignment refines transcript timestamps so text aligns with
audio/video at word level \cite{mcauliffe2017montreal}.
\begin{itemize}
    \item \textbf{Core technology:} Forced Alignment algorithms, such as
          Wav2Vec2 \cite{baevski2020wav2vec} or WhisperX
          \cite{bain2023whisperx}, synchronize each word with the exact
          millisecond at which it is spoken in the original video.
    \item \textbf{Role in LMS:} This process is crucial for real-time
          temporal context. It allows the system to determine exactly when
          a concept is taught, enabling the LMS to pause the video at the
          right moment for Pop-up Quiz insertion or interactive
          transcript features.
\end{itemize}

\subsection{Video Segmentation}
Video Segmentation divides a long lecture video into shorter
micro-learning clips, each focused on a specific academic topic or
concept \cite{galanopoulos2020lecture}.
\begin{itemize}
    \item \textbf{Mechanism:} Segmentation combines several analysis streams:
    \begin{itemize}
        \item \textit{Semantic Text Segmentation:} uses LLMs to analyze
              transcript semantics and identify topic shifts.
        \item \textit{Visual Cues:} applies Computer Vision to detect
              slide changes, scene cuts, or board interactions.
        \item \textit{Audio Pauses:} identifies intentionally long pauses
              in the instructor's speech.
    \end{itemize}
    \item \textbf{Role in LMS:} The output is independent micro-content
          video segments. Each segment is then automatically summarized,
          tagged, and titled by AI to optimize personalized learning.
\end{itemize}

\subsection{Multimodal Learning Pipeline}
The Multimodal Learning Pipeline integrates multiple data streams to
provide the richest context for Generative AI \cite{baltrusaitis2018multimodal}.
\begin{itemize}
    \item \textbf{Overall architecture:}
    \begin{enumerate}
        \item \textbf{Audio stream:} data passes through Whisper and
              Transcript Alignment to create word-level timestamped text.
        \item \textbf{Video-frame stream:} data passes through Computer
              Vision and OCR to extract text from lecture slides and
              record visual transition events.
        \item \textbf{Multimodal fusion:} spoken transcript and slide text
              are tightly synchronized on the same timeline to form a
              shared metadata base.
        \item \textbf{Generative AI processing:} the multimodal metadata
              block is passed into an LLM (with multimodal support) along
              with Prompt Engineering instructions to automatically
              generate quiz questions and segment videos logically.
    \end{enumerate}
    \item \textbf{Role in LMS:} Multimodal fusion prevents the LLM from
          missing important information when instructors explain using
          visual elements, improving the accuracy of micro-content and
          quiz context.
\end{itemize}

\section{AI Orchestration Frameworks and Supporting Libraries}
In complex Generative AI applications, a Large Language Model rarely works
alone. To build a multimodal pipeline from video cutting and transcript
reading to quiz generation, the system needs an ``orchestrator'' to
coordinate data, manage memory, and connect tasks. AI Orchestration
Frameworks solve this problem.

\subsection{LangChain}
LangChain is a leading open-source platform designed to simplify the
development of LLM-based applications \cite{chase2022langchain}. In the
LMS system, LangChain acts as a framework connecting models such as
Whisper or GPT-4 with lecture data and internal databases. Core LangChain
components applied in the system include:

\begin{itemize}
    \item[\textbf{a.}] \textbf{Document Loader:} \\
    This is the first step in the data-processing flow. Document Loader
    provides standard interfaces to extract data from many formats and
    sources, such as TXT, JSON, PDF lecture slides, or video transcripts.
    In the LMS context, loaders take raw data from speech recognition
    (Whisper) and OCR, then normalize them into \texttt{Document} objects
    containing text and metadata such as timestamps for later analysis.

    \item[\textbf{b.}] \textbf{Retriever:} \\
    This component is the heart of Retrieval-Augmented Generation
    \cite{lewis2020retrieval}. When the LMS needs to create a quiz for a
    video segment from minute 05:00 to 10:00, the Retriever finds and
    extracts the correct transcript passages, slide content, and key
    concepts for that time range from the Vector Store. This gives the LLM
    accurate context and prevents hallucination during question generation.

    \item[\textbf{c.}] \textbf{Prompt Template:} \\
    Prompt Template parameterizes and reuses prompts sent to the LLM
    \cite{white2023prompt}. Instead of rewriting prompts from scratch for
    each video segment, the system uses dynamic templates. For example:
    \textit{"Based on the lecture transcript: \{transcript\_segment\} and
    board content: \{ocr\_data\}, create \{num\_questions\}
    multiple-choice questions at difficulty \{difficulty\_level\}."}
    LangChain injects these variables into the prompt before calling the
    language model API.

    \item[\textbf{d.}] \textbf{Chain Orchestration:} \\
    This mechanism links multiple sequential or parallel processing steps
    into a complete pipeline \cite{wu2022promptchainer}. In the
    micro-content generation system, a typical LLMChain proceeds as:
    (1) load data from Document Loader $\rightarrow$
    (2) Retriever fetches relevant context $\rightarrow$
    (3) pass data into Prompt Template $\rightarrow$
    (4) send to the LLM for question generation $\rightarrow$
    (5) pass through an Output Parser to coerce output into a standard
    format such as JSON before saving it to the LMS database. Chains allow
    the process to be automated without manual intervention between steps.
\end{itemize}

\subsection{LangGraph}
Although LangChain provides powerful tools for connecting tasks (Chains),
those chains are usually linear and less flexible for complex scenarios
requiring loops. LangGraph extends the LangChain ecosystem by allowing LLM
applications to be built as cyclic graphs with complex state management
\cite{langgraph2024}. In the LMS system, LangGraph provides
infrastructure for reliable automated content generation through three
core properties:

\begin{itemize}
    \item[\textbf{a.}] \textbf{Stateful Workflow:} \\
    Unlike stateless chains, LangGraph maintains a Global State throughout
    execution. In the LMS context, this state may store the original
    video, the transcript being processed, extracted key concepts, and the
    generated quiz questions. As data passes through nodes in the graph,
    for example from video cutting to question generation, this state is
    continuously updated and inherited, so the system does not lose
    context in long multi-step workflows.

    \item[\textbf{b.}] \textbf{Agentic Pipeline:} \\
    LangGraph allows the system to be designed as multiple autonomous AI
    Agents that coordinate with each other \cite{wang2023survey}. Instead
    of using one LLM for everything, the LMS can split tasks into: (1)
    \textit{Extraction Agent} for semantic analysis and video cutting, (2)
    \textit{Quiz Generation Agent} for multiple-choice question
    generation, and (3) \textit{Reviewer Agent} for quality control. These
    agents can call tools and communicate, creating a specialized and more
    optimized processing pipeline.

    \item[\textbf{c.}] \textbf{Branching Flow:} \\
    The major strength of graphs is conditional edges. Instead of always
    moving from A to B, the system can branch or loop based on evaluation
    results \cite{wei2022chain}. For example, after the Quiz Generation
    Agent creates a question, the workflow moves to review. If the
    question lacks an answer or does not follow the source transcript, the
    workflow loops back to the generation node with an error message for
    the LLM to repair. This self-correction branching helps ensure that
    micro-content and quizzes satisfy pedagogical quality before entering
    the LMS database.
\end{itemize}

\subsection{Document Processing Libraries}
Besides audio and visual data from video, the LMS system often needs to
process attached documents such as lecture slides (PowerPoint), textbooks
(PDF), or text documents (Word). Document Processing Libraries extract raw
text, tables, and metadata from these files, converting them into a
structured format for the RAG system to generate quiz questions.

\begin{itemize}
    \item[\textbf{a.}] \textbf{Docling:} \\
    Docling is an advanced open-source library developed by IBM for
    parsing and understanding complex documents \cite{docling2024}. In an
    LMS, Docling is especially useful for PDFs containing mathematical
    formulas, scientific tables, and diagrams. It can extract and preserve
    document structure in Markdown or JSON, helping the LLM maintain
    context when reading complex academic content.

    \item[\textbf{b.}] \textbf{Unstructured:} \\
    Unstructured is a standard library in the RAG ecosystem for extracting
    unstructured data \cite{unstructured2024}. It provides unified APIs to
    read dozens of formats such as HTML, PDF, TXT, and EPUB. In the
    micro-content generation system, Unstructured works well with
    LangChain for Document Chunking, splitting textbooks into meaningful
    short passages so the Vector Database can retrieve them easily.

    \item[\textbf{c.}] \textbf{PyMuPDF:} \\
    PyMuPDF is a high-performance PDF processing library based on the
    MuPDF engine \cite{mckennon2024pymupdf}. Its strength in the LMS
    pipeline is very fast text extraction and direct rendering of PDF
    pages into image frames. The system can use PyMuPDF to extract text
    from PDF slides and align it with video display time.

    \item[\textbf{d.}] \textbf{python-docx:} \\
    This is the standard Python library for creating and updating
    Microsoft Word (.docx) files \cite{python_docx2024}. In this system,
    it has two purposes: (1) read and extract scripts or lesson plans
    uploaded by instructors; (2) export generated results, for example a
    set of 50 multiple-choice questions with answers, into a formatted
    Word file so instructors can download and manually edit them if
    needed.

    \item[\textbf{e.}] \textbf{PPTX Parser (python-pptx):} \\
    Lecture slides are a core knowledge source. \texttt{python-pptx}
    allows the system to inspect .pptx files and extract text from each
    Text Box and Speaker Notes on each slide \cite{python_pptx2024}. This
    data is then aligned with the video transcript to create a highly
    accurate multimodal learning context.
\end{itemize}

\subsection{Video Processing Libraries}
After AI Agents such as LLMs analyze semantics and decide the video cut
points, for example cutting from second 120 to 340 to form micro-content
lesson 1, the system needs low-level tools to execute heavy media-file
processing.

\begin{itemize}
    \item[\textbf{a.}] \textbf{FFmpeg:} \\
    FFmpeg is the ``gold standard'' in multimedia processing
    \cite{tomar2006converting}. In the LMS system, FFmpeg runs at the
    operating-system level and performs the heaviest tasks at high speed
    with efficient resource usage: extracting the audio stream (.wav) from
    the original lecture video (.mp4) for speech recognition, or extracting
    batches of frames for Computer Vision analysis.

    \item[\textbf{b.}] \textbf{MoviePy:} \\
    MoviePy is a Python library providing a high-level, easy-to-use API
    around the FFmpeg engine \cite{moviepy2024}. The system uses MoviePy
    to execute automatic subclip cutting, video concatenation, fade-in and
    fade-out effects, and insertion of text or watermarks such as the
    university logo or micro-lesson title over the output video before
    storing it in LMS storage.

    \item[\textbf{c.}] \textbf{Whisper (Library Integration):} \\
    Besides being an AI model, Whisper is also available as the open-source
    Python library \texttt{openai-whisper} \cite{radford2023robust}. This
    library integrates closely with PyTorch and FFmpeg, allowing the system
    to programmatically load models to GPU, set decoding parameters,
    perform language detection, and return JSON objects containing
    transcripts with detailed word-level timestamps.
\end{itemize}

\section{Technology Selection Rationale}
\label{sec:technology-selection-rationale}

This section does not repeat textbook definitions of common tools. The
important point for this thesis is why each technology fits the proposed
AI-enabled LMS architecture. The theoretical core of the chapter remains
RAG, GraphRAG, chunking, Bloom-aware quiz generation, and LTI 1.3; the
following technologies are implementation choices that support those
foundations.

\begin{itemize}
    \item \textbf{FastAPI} \cite{ramirez2018fastapi}: selected for the Python
          AI backend because it integrates naturally with parsing, embedding,
          LLM, and video-processing libraries while supporting asynchronous
          API calls.
    \item \textbf{Go} \cite{golang2026}: selected for the video worker because
          video processing requires robust subprocess orchestration, concurrent
          task execution, and predictable resource usage around FFmpeg,
          transcription, and media-cutting jobs.
    \item \textbf{React + Next.js} \cite{vercel2024nextjs}: selected for
          instructor and learner interfaces because the system needs
          dashboard screens, review editors, quiz views, and LMS entry routes.
    \item \textbf{PostgreSQL} \cite{postgresql2024}: selected as the main
          relational store for users, courses, documents, attempts, review
          states, and transactional learning data.
    \item \textbf{Redis} \cite{carlson2013redis}: selected as the queue/cache
          layer so upload, transcription, embedding, and LLM generation jobs
          can run asynchronously.
    \item \textbf{Qdrant} \cite{qdrant2024}: selected as the vector database
          for RAG because it supports embedding search with metadata filters.
    \item \textbf{Neo4j} \cite{neo4j2024}: selected for curriculum and
          Learning Outcome relationships that support GraphRAG expansion.
    \item \textbf{MinIO} \cite{minio2024}: selected for large media objects
          such as lecture videos, extracted audio, rendered pages, and
          generated files instead of storing them in PostgreSQL.
    \item \textbf{gRPC} \cite{grpc2024}: selected for internal service calls
          where typed contracts and efficient binary transport are useful
          between web, core, and AI worker runtimes.
    \item \textbf{Model Context Protocol (MCP)}
          \cite{modelcontextprotocol2025}: selected for external AI tool
          integration, especially the YouTube/video-search adapter, because it
          standardizes how the system exposes tools and keeps third-party API
          access behind a controlled adapter boundary.
    \item \textbf{Docker} \cite{merkel2014docker}: selected to package the
          mixed runtime stack consistently across local and deployment
          environments.
\end{itemize}

These choices are not the research contribution by themselves. Their role
is to provide a practical execution environment for the core ideas: source
grounding through RAG, structural retrieval through GraphRAG, pedagogical
alignment through Bloom and Learning Outcomes, and LMS interoperability
through LTI 1.3.
